\documentclass[tikz,border=5pt]{standalone}
\usepackage{luatexja}
\usepackage[hiragino-pron]{luatexja-preset}
\usepackage{amsmath,amssymb}
\usetikzlibrary{arrows.meta,patterns,calc}

\begin{document}
\begin{tikzpicture}

% ============================================================
% Panel 1 (Left): 2D domain D with horizontal slice at y
% ============================================================
\begin{scope}[shift={(0,0)}]

% Axes
\draw[-{Stealth[length=6pt]}, thick] (-0.5,0) -- (5,0) node[right] {$x$};
\draw[-{Stealth[length=6pt]}, thick] (0,-0.5) -- (0,4) node[above] {$y$};

% Region D (ellipse)
\draw[thick] (2.4,2.0) ellipse [x radius=1.6, y radius=1.0];

% D label
\node at (2.4,3.4) {$D$};

% Horizontal blue line at y (slice)
% Ellipse at (2.4,2.0), rx=1.6, ry=1.0. At y=2.3: dx = 1.6*sqrt(1-((2.3-2.0)/1.0)^2) = 1.6*0.954 = 1.526
\draw[blue, very thick] ({2.4-1.526},2.3) -- ({2.4+1.526},2.3);

% Dashed lines to axes
\draw[dashed, thin] (0,1.0) -- (2.4,1.0);
\draw[dashed, thin] (0,3.0) -- (2.4,3.0);
\draw[dashed, thin] ({2.4-1.526},0) -- ({2.4-1.526},2.3);
\draw[dashed, thin] ({2.4+1.526},0) -- ({2.4+1.526},2.3);
\draw[blue, dashed, thin] (0,2.3) -- ({2.4-1.526},2.3);

% Labels
\node[left] at (0,1.0) {$y_0$};
\node[left] at (0,3.0) {$y_1$};
\node[below] at ({2.4-1.526},0) {$c(y)$};
\node[below] at ({2.4+1.526},0) {$d(y)$};

\end{scope}

% ============================================================
% Panel 2 (Middle): 3D volume (cylinder with slice at fixed y)
% ============================================================
\begin{scope}[shift={(9.5,0)}]

\coordinate (O) at (-1,0.8);

% Axes
\draw[-{Stealth[length=6pt]}, thick] (O) -- ++(0,4.0) node[above] {$z$};
% y-axis: solid before cylinder, dashed behind, solid after
\draw[thick] (O) -- ++(-0.5+1.0,0);          % origin to left edge of cylinder
\draw[thick, dashed] (-0.5,0.8) -- (2.1,0.8); % behind cylinder
\draw[-{Stealth[length=6pt]}, thick] (2.1,0.8) -- (4.0,0.8) node[right] {$y$};
\draw[-{Stealth[length=6pt]}, thick] (O) -- ++(-2.0,-1.2) node[below left] {$x$};

% Domain D on xy-plane (ellipse, no hatching)
\coordinate (Dc) at (0.8,0.05);
\begin{scope}[shift={(0.8,0.05)}, rotate=-5]
  % Back half (dashed)
  \draw[thick, dashed] (180:1.3 and 0.5) arc[x radius=1.3, y radius=0.5, start angle=180, end angle=0];
  % Front half (solid)
  \draw[thick] (180:1.3 and 0.5) arc[x radius=1.3, y radius=0.5, start angle=180, end angle=360];
\end{scope}

% D label
\node at ($(Dc)+(1.7,-0.2)$) {$D$};

% Curve on surface (top of cylinder)
\draw[thick]
  (-0.5,2.65) .. controls (-0.3,3.35) and (0.5,3.55) .. (1.0,3.35)
  .. controls (1.5,3.15) and (2.2,3.05) .. (2.1,2.65)
  .. controls (2.0,2.25) and (1.3,2.15) .. (0.8,2.25)
  .. controls (0.2,2.35) and (-0.6,2.25) .. (-0.5,2.65);

% Vertical lines (solid)
\draw[thick] (-0.5,2.60) -- (-0.5,0.15);
\draw[thick] (2.1,2.80) -- (2.1,0.04);

% Cross-section T(y) - red shaded region (parallel to xz-plane)
% Base runs along x-direction (-1,-0.6), sides go straight up (z)
% Base endpoints lie ON the domain ellipse
% Ellipse: center (0.8,0.05), rx=1.3, ry=0.5 (approx, ignoring -5deg rotation)
% Chord with slope 0.6 (= x-direction): y = 0.6x - 0.4
% Solving with ellipse gives: (1.47, 0.48) and (0.06, -0.36)
\coordinate (base-c) at (1.47, 0.48);   % c(y): near origin on ellipse
\coordinate (base-d) at (0.2, -0.35);  % d(y): far along x on ellipse
\coordinate (top-c) at (1.47, 3.20);    % straight up from base-c
\coordinate (top-d) at (0.2, 2.30);    % straight up from base-d

\fill[red, opacity=0.15]
  (base-d) -- (top-d)
  .. controls (0.5,2.10) and (1.0,2.80) .. (top-c)
  -- (base-c) -- cycle;
\draw[red, thick] (base-d) -- (top-d);
\draw[red, thick] (top-d) .. controls (0.5,2.10) and (1.0,2.80) .. (top-c);
\draw[red, thick] (base-c) -- (top-c);
\draw[red, thick] (base-d) -- (base-c);

% Dashed blue lines projecting c(y), d(y) onto x-axis
% Project horizontally (y-direction) from base endpoints to x-axis line
% x-axis: O=(-1,0.8) + t*(-1,-0.6) → (fig_x, fig_y) = (-1-t, 0.8-0.6t)
% base-c at fig_y=0.48: t=(0.8-0.48)/0.6=0.533, fig_x=-1.53
% base-d at fig_y=-0.36: t=(0.8+0.36)/0.6=1.933, fig_x=-2.93
\coordinate (cy-on-axis) at (-1.53, 0.48);
\coordinate (dy-on-axis) at (-2.92, -0.35);
\draw[blue, dashed, thick] (base-c) -- (cy-on-axis);
\draw[blue, dashed, thick] (base-d) -- (dy-on-axis);
\node[blue, above left] at (cy-on-axis) {$c(y)$};
\node[blue, above left] at (dy-on-axis) {$d(y)$};

% T(y) label
\node[red] at (0.8,1.6) {$T(y)$};

\end{scope}

% ============================================================
% Panel 3 (Right): Cross-section T(y) = integral over x
% ============================================================
\begin{scope}[shift={(15,0.5)}]

% Axes
\draw[-{Stealth[length=6pt]}, thick] (-0.5,0) -- (4,0) node[right] {$x$};
\draw[-{Stealth[length=6pt]}, thick] (0,-0.5) -- (0,3.5) node[above] {$z$};

% f(x,y) curve (cross section at fixed y)
\draw[thick, domain=0.3:3.5, samples=60, smooth]
  plot (\x, {0.35*(\x-1.8)*(\x-1.8)+1.8});

% Shaded area T(y)
\fill[blue, opacity=0.15]
  (0.5,0) -- plot[domain=0.5:3.1, samples=60, smooth]
  (\x, {0.35*(\x-1.8)*(\x-1.8)+1.8}) -- (3.1,0) -- cycle;

% Boundary marks
\draw[dashed, thin] (0.5,0) -- (0.5,{-0.35*(0.5-1.8)*(0.5-1.8)+2.5});
\draw[dashed, thin] (3.1,0) -- (3.1,{-0.35*(3.1-1.8)*(3.1-1.8)+2.5});

% Labels
\node[below] at (0.5,0) {$c(y)$};
\node[below] at (3.1,0) {$d(y)$};
\node[above right] at (3.3,1.8) {$f(x,y)\ \text{（$y$固定）}$};

% T(y) label with integral
\node[blue!70!black, align=center, inner sep=0pt] at (1.8,1.0) {$T(y)$\\[4pt]{\scriptsize $\displaystyle= \int_{c(y)}^{d(y)} f(x,y)\,dx$}};

\end{scope}

\end{tikzpicture}
\end{document}
